\section{Removing multiplicative masks and detecting a shifted error}
\label{sec:correlation}

The purpose of this Section is to turn a correlation containing arbitrary
multiplicative masks into a test of one function on an additive cube.
The number of Cauchy--Schwarz steps will depend only on the number of
blocks in the chain.  This independence from the master index count is
needed when the analytic tolerances are chosen before the Ramsey bound.

Fix a master chain
\[
 B_k=T_k\cup\{i_k\},\qquad
 T_1<\cdots<T_m<l<i_1<\cdots<i_m,
\]
where each tail is nonempty, and fix positive scales
$c_k=h_{B_k}a_k$ from the finite list allowed in
Section~\ref{sec:arithmetic}.  Write $V=V_l$, $R=R_l$, and let
$\mathcal P_l$ be the prime pool in this gap, with law
\[
 \lambda(p)=\frac{1}{pS_l},\qquad
 S_l=\sum_{p\in\mathcal P_l}\frac1p.
\]
In this Section $z_k$ has law $\mu_{i_k}$, independently for different
$k$.  We abbreviate $\nu_k=\nu_{B_k}$, so that $0\le\nu_k\le V$ on
all integers.  For $\emptyset\ne J\subseteq[m]$ put
\[
 a(J)=\max J,\qquad
 L_J(z)=\sum_{k\in J}\frac{c_k}{c_{a(J)}}z_k.
\]
For sufficiently large $w$, all the displayed coefficients are integers,
and every coefficient other than the anchor coefficient is divisible by
$W$.  The scales $c_k$ are positive integers with no prime factor greater
than $w$.  We consider
\begin{equation}
 \mathcal C=
 \E_z\left[
   \prod_{\emptyset\ne U\subseteq[m]}b_U(z_U)
   \prod_{\emptyset\ne J\subseteq[m]}g_J(L_J(z))
 \right],
 \qquad z_U=\prod_{k\in U}z_k,
 \label{eq:correlation-initial}
\end{equation}
where the functions are real, $|b_U|\le1$, and
$|g_J(y)|\le1+\nu_{a(J)}(y)$ for every integer $y$.
Fix $J_*$ with $|J_*|\ge2$, and write $g_*=g_{J_*}$ and
$a_* = a(J_*)$.  Bounded extensions outside the positive integers are
permitted whenever needed.

We will keep one copy of $g_*$ throughout: first remove the product
masks, then eliminate the other linear factors, and finally identify
the sampling law of the surviving cube root.

All estimates below are uniform over these functions.  An $o(1)$ may
depend on the fixed master data and on a fixed positive integer $J_0$,
but not on the functions.  Uniformity also holds over any fixed finite
set of $J_0$'s.  In particular, it does not assert a convergence rate
uniform in a growing master index count.

\subsection{Prime insertion and weighted mask removal}

\begin{lemma}[Prime insertion]\label{lem:prime-insertion}
Let $i>l$ and let $Y$ have law $\mu_i$.  For every fixed $A>0$, uniformly
over functions $F$ with $|F|\le V^A$,
\begin{equation}
 \E_Y F(Y)=\E_{p,Y}F(pY)+o(1).
 \label{eq:prime-average-insertion}
\end{equation}
For a fixed positive integer $k$, coprime to $W$ and bounded by a fixed
power of $P_l^++V$, one also has
\begin{equation}
 \E_Y F(kY)=\E_Y k\1_{k\mid Y}F(Y)+o(1),
 \label{eq:prime-fixed-dilation}
\end{equation}
uniformly in $k$.  The total-mass error before multiplication by $F$ is
smaller than every fixed negative power of $P_l^++V$.
These assertions remain valid with other independent variables as
parameters.
\end{lemma}

\begin{proof}
The dilation assertion is Lemma~\ref{lem:sampling}, applied at a cutoff
whose logarithm dominates all fixed powers of $P_l^++V$.  The weighted
law on the right has the same harmonic density as the law on the left;
only the endpoints $[X_i,X_i^2)$ and $[kX_i,kX_i^2)$ differ.
The logarithmic boundary mass is $O(\log k/\log X_i)$, with the residue
counting error from that Lemma.  The scale separation makes their sum
smaller than any prescribed negative power of $P_l^++V$.

Averaging the multiplier in \eqref{eq:prime-fixed-dilation} gives
\[
 A_l(Y)=\E_p p\1_{p\mid Y}
       =\frac1{S_l}\sum_{p\in\mathcal P_l}\1_{p\mid Y}.
\]
Periodic counting for the moduli $p$ and $pq$, including the $W$-unit
restriction, gives their reciprocal probabilities with uniformly
negligible relative errors.  Therefore
\begin{align*}
 \E A_l&=1+o(V^{-C}),\\
 \E A_l^2
 &=\frac1{S_l^2}
   \left(\sum_p\frac1p+\sum_{p\ne q}\frac1{pq}\right)
   +o(V^{-C})\\
 &=1+\frac1{S_l}-\frac1{S_l^2}\sum_p\frac1{p^2}+o(V^{-C})
\end{align*}
for every fixed $C$.  Thus
$\|A_l-1\|_{L^2(\mu_i)}\le S_l^{-1/2}+o(V^{-C})$.
Cauchy--Schwarz bounds the error in replacing $A_l$ by $1$ by
$V^A(S_l^{-1/2}+o(V^{-C}))=o(1)$, since $S_l$ dominates every fixed
power of $V$.  The estimates are uniform in the parameters of $F$, so
one may integrate over those parameters afterwards.
\end{proof}

A row template will mean a vector $A_R=(A_{R,1},\ldots,A_{R,m})$ whose
nonzero entries are monomials in finitely many formal prime variables.
No prime variable occurs in two different columns of a single row.
For $a(R)=\max\{k:A_{R,k}\ne0\}$, its associated linear form and weight
are
\begin{equation}
 \ell_R(z)=\sum_k\frac{c_k}{c_{a(R)}}A_{R,k}z_k,
 \qquad W_R(y)=1+\nu_{a(R)}(y).
 \label{eq:correlation-row-template}
\end{equation}
Parallelity of templates always means parallelity over the rational
function field in the formal prime variables.  This qualification is
essential: accidental equalities after numerical substitution are
exceptional events, not identities of templates.

Reciprocal dilations preserving a product, followed by
Cauchy--Schwarz to remove its bounded factor, also occur in the proof
of \cite[Proposition~2.4]{GreenSanders} over finite fields.
Here the divisibility weights allow the corresponding substitutions
on integer variables, and we carry out the elimination for the full
family of product masks.

\begin{lemma}[Weighted removal of multiplicative masks]
\label{lem:mask-removal}
Put $q_{\mathrm{mask}}=2^m-1$ and $K_m=q_{\mathrm{mask}}2^{q_{\mathrm{mask}}}$.  Starting from
\eqref{eq:correlation-initial}, one obtains a family $\mathcal R$ of at
most $K_m$ pairwise nonparallel row templates, with one distinguished
row $*$, and functions $f_R$ that may depend on the introduced primes,
such that
\begin{equation}
 |\mathcal C|^{2^{q_{\mathrm{mask}}}}
 \le C_m\left|\E_{\boldsymbol p,z}
                 \prod_{R\in\mathcal R}f_R(\ell_R(z))\right|+o(1).
 \label{eq:mask-removal-output}
\end{equation}
Here $|f_R|\le W_R$, the distinguished function is exactly $f_*=g_*$,
its support is $J_*$, and every prime parameter has the law $\lambda$,
independently before deletion of exceptional tuples.
The row templates, the number of parameters, and $C_m$ depend only on
$m$ and the choice of $J_*$.  They are independent of the scales and
functions.
\end{lemma}

\begin{proof}
Initially the rows are the indicator vectors of the nonempty subsets
$J$ of $[m]$.  They are pairwise nonparallel.  We remove the masks in a
fixed order, maintaining \eqref{eq:correlation-row-template}, the weight
bounds, and a distinguished copy of $g_*$.

\paragraph{A substitution fixing the mask.}
Consider the mask with support $U$.  If the distinguished support has
an index $u\notin U$, use the change $D(p)z$ given by
$z_u\mapsto pz_u$.  Otherwise choose distinct $u,v$ in the distinguished
support, necessarily both in $U$, and use
\begin{equation}
 z_u\mapsto z_u/p,\qquad z_v\mapsto pz_v,
 \qquad \eta_p(z)=p\1_{p\mid z_u}.
 \label{eq:balanced-prime-substitution}
\end{equation}
The first case has $\eta_p=1$.
Lemma~\ref{lem:prime-insertion} justifies the first case directly.
For the second, first insert $p$ in coordinate $v$ and then use
\eqref{eq:prime-fixed-dilation} in coordinate $u$ with the function
read at $z_u/p$.  Thus division by $p$ is used only on the support of
the displayed indicator, and its multiplier is retained.  At this
stage all unweighted functions and all products of row weights are
bounded by a fixed power of $V$; even the intermediate fixed-prime
multipliers are absorbed by the stronger endpoint error in
\eqref{eq:prime-fixed-dilation}.  Both substitutions therefore have
$o(1)$ error.  In either case $z_U$ is unchanged.

Call a row invariant at this step when the coefficient vector of
$\ell_R(D(p)z)$ is a scalar multiple of that of $\ell_R(z)$.
The scalar is $1,p$, or $p^{-1}$.  All divisors in $\nu_{a(R)}$ are
smaller than the pool primes, so multiplication or admissible division
by $p$ preserves their divisibility conditions.  Consequently
\[
 W_R(\ell_R(D(p)z))=W_R(\ell_R(z))
\]
on the support in use.  Put
$\Omega(z)=\prod_{R\text{ invariant}}W_R(\ell_R(z))$.
Divide each invariant function inside the fresh prime average by its
corresponding weight, and call the remaining integrand $H_p$.
Let $\boldsymbol r$ denote all prime slots introduced at earlier steps.
The current mask and row functions may depend on $\boldsymbol r$.
All these previous slots are integrated together with $z$ in the outer
expectation. Thus the substituted correlation is
\[
 I=\E_{\boldsymbol r,z}
 b_U(\boldsymbol r;z_U)\Omega(\boldsymbol r,z)
       \E_p\eta_p(z)H_p(\boldsymbol r,z).
\]
Here the new slot $p$ is independent of $(\boldsymbol r,z)$.
Weighted Cauchy--Schwarz on the joint outer probability space gives
\begin{equation}
 |I|^2\le
 \bigl(\E_{\boldsymbol r,z}\Omega(\boldsymbol r,z)\bigr)
 \E_{\boldsymbol r,z,p,q}
 \Omega(\boldsymbol r,z)\eta_p(z)\eta_q(z)
 H_p(\boldsymbol r,z)H_q(\boldsymbol r,z).
 \label{eq:mask-weighted-cs}
\end{equation}
The old tuple $\boldsymbol r$ is shared by the two branches, while $p$
and $q$ are fresh independent slots. This inequality applies to the
full prime average.

If there are $s_{\rm inv}$ invariant rows, the first factor equals
$2^{s_{\rm inv}}+o(1)$, and hence is bounded in terms of the number of
rows alone. Indeed, expand $\Omega$ as a sum of products of $\nu$'s
and apply Proposition~\ref{prop:linear-forms} with $\boldsymbol r$
among its prime parameters. Every anchor coefficient is a monomial
of pool primes, hence a unit at every prime $w<\pi\le V$. A nonzero
minor of two templates is multiplied in the actual rows by smooth
row and column factors, also units at such $\pi$, so its possible
vanishing is covered by a fixed polynomial test. This bound uses
the average over $\boldsymbol r$.

\paragraph{The two prime branches.}
In the one-coordinate case the two branches have multipliers $p$ and
$q$ on $u$.  In the two-coordinate case, first discard $p=q$.
Since all remaining factors are bounded by $V^{O_m(1)}$, its contribution
is at most
\[
 V^{O_m(1)}\sum_p\lambda(p)^2p^2\PP(p\mid z_u)
 \le (1+o(1))V^{O_m(1)}\sum_p p\lambda(p)^2
 =\frac{V^{O_m(1)}}{S_l}+o(1)=o(1).
\]
For $p\ne q$, the multiplier is $pq\1_{pq\mid z_u}$.
Absorb it by \eqref{eq:prime-fixed-dilation}, replacing $z_u$ by
$pqz_u$.  The two branch multipliers on $(u,v)$ are then
\begin{equation}
 (q,p)\quad\hbox{and}\quad(p,q).
 \label{eq:prime-two-branches}
\end{equation}
Every surviving multiplicative mask is a bounded function of a scalar
multiple of its original monomial; its two copies can be combined into
one mask on that same support.

For an invariant row the two scalar-changed arguments are
$\alpha_p\ell_R(z)$ and $\alpha_q\ell_R(z)$, after absorption if
necessary.  Its new function is
\[
 t\longmapsto
 \frac{f_R(\alpha_p t)f_R(\alpha_q t)}{W_R(t)},
\]
on the arguments actually used, and it is bounded by $W_R(t)$.
This formula records why only one weight survives.  It is extended
elsewhere with the same bound.  In particular, in the two-coordinate
case an invariant singleton on $u$ acquires arguments $qt$ and $pt$;
the weight originally evaluated at $pqt$ equals $W_R(t)$.
Every noninvariant row instead gives two functions, each with its
original weight bound, on the two substituted rows.

\paragraph{Preserving the row structure.}
The new templates are pairwise nonparallel.  For descendants of
different old rows, a purported symbolic parallelity specializes at
$p=q=1$ to parallelity of the old rows, a contradiction.  For the two
descendants of one row, use the exponent description of the
substitution: a column is multiplied by $p^{e_k}$, with
$e_k\in\{0,1\}$ or $\{-1,0,1\}$.  The $p$ and $q$ branches are
parallel precisely when all $e_k$ on the row support are equal.
That is exactly the definition of an invariant row.  The common
column multiplication by $pq$ in \eqref{eq:prime-two-branches} is
invertible over the rational function field and does not affect this
argument.  The distinguished row is noninvariant: it contains both an
affected and a differently affected coordinate.  Track a fixed one
of its two copies.  Its function is still $g_*$ and its support is
unchanged.

The new slots occur in at most one column per row, as is explicit in
\eqref{eq:prime-two-branches}.  Thus the template invariant persists.
Repeated prime entries and exact polynomial coincidences have
probability smaller than every fixed negative power of $V$.
After absorption the integrands are bounded by $V^{O_m(1)}$, so these
events can be deleted or reinstated at $o(1)$ cost.  The reinstated
substitution-form terms on $p=q$ have this same harmless bound; the
larger, pre-absorption diagonal was estimated separately above.
Consequently fresh slots can always be sampled independently.
At the next step the entire existing tuple, including this step's
$p$ and $q$, is the shared outer tuple $\boldsymbol r$. The deletion
and reinstatement argument restores its independent product law
before that next step. Any temporarily retained prime-only domain
is covered by the restricted form of
Proposition~\ref{prop:linear-forms}; the base variables and the
independently expanded divisor draws retain their original laws.

Each step removes one mask, uses one square, and at most doubles the
number of rows.  There are $q_{\mathrm{mask}}$ masks and initially $q_{\mathrm{mask}}$ rows, giving
at most $K_m$ rows.  Iterating \eqref{eq:mask-weighted-cs}, with its
bounded prefactors and uniform errors, proves
\eqref{eq:mask-removal-output}.  Fixing the order of masks, the choices
of coordinates, and the tracked branch makes every template depend
only on $m$ and $J_*$.  The original and intermediate correlations are
bounded by constants, by the same linear-forms estimates, so raising
the inequalities to these fixed powers preserves $o(1)$ errors.
\end{proof}

\subsection{Directions with one vanishing row response}

The product masks have been removed.  To eliminate a nontarget row by
Cauchy--Schwarz, we need a translation that fixes that row and moves
every other row, including the target.  We choose these translations so
that all target responses are equal.

Let $\mathcal R$ now denote the output family of
Lemma~\ref{lem:mask-removal}, and let $d=|\mathcal R|-1$.
For a nonzero integer $n$, write $|n|_{>w}$ and $|n|_{\le w}$ for its
rough and smooth parts, including multiplicity.

\begin{lemma}[Polynomial directions and integer translations]
\label{lem:row-directions}
For each $R\ne *$ there is an integer polynomial vector $w_R$ such that
\[
 A_Rw_R=0,\qquad A_Iw_R\ne0\quad(I\ne R).
\]
There is also an integer polynomial vector $w_0$ with
$A_*w_0=0$ and $A_Iw_0\ne0$ for $I\ne *$.
These choices can be made from a finite list determined solely by the
row templates.  Set
\begin{equation}
 d_R=A_*w_R,\qquad D=\prod_{R\ne *}d_R,\qquad
 M(\boldsymbol p)=M|D(\boldsymbol p)|_{>w}.
 \label{eq:correlation-modulus}
\end{equation}
On the good prime tuples planned in Lemma~\ref{lem:master-scales},
the vectors
\begin{equation}
 v_{R,k}=M(\boldsymbol p)\frac{c_{a_*}}{c_k}
                  \frac{w_{R,k}}{d_R},\qquad
 v_{0,k}=M\frac{c_{a_*}}{c_k}w_{0,k}
 \label{eq:correlation-integer-directions}
\end{equation}
belong to $W\Z^m$ and have sizes bounded by a fixed power of
$P_l^++V$.  Their responses satisfy
\begin{equation}
 \ell_*(v_R)=M(\boldsymbol p),\qquad
 \ell_R(v_R)=0,\qquad \ell_*(v_0)=0.
 \label{eq:correlation-direction-responses}
\end{equation}
Every other response is nonzero.  At primes $w<\pi\le V$ it is a unit
unless a member of a fixed list of nonzero integer polynomials in the
prime slots vanishes modulo $\pi$.
\end{lemma}

\begin{proof}
Here is an explicit finite choice of a kernel vector.  For a row $A_R$,
choose a column $j$ with $A_{R,j}\ne0$.  The polynomial vectors
\[
 b_k=A_{R,j}e_k-A_{R,k}e_j\qquad(k\ne j)
\]
span its kernel over the rational function field.  Enumerate them as
$b_0,\ldots,b_{m-2}$ and set $w_R(t)=\sum_{h=0}^{m-2}t^hb_h$.
For $I\ne R$, the polynomial $A_Iw_R(t)$ is not identically zero:
otherwise $A_I$ would vanish on the kernel of $A_R$ and hence be
parallel to it.  The product of these polynomials has degree at most
$(|\mathcal R|-1)(m-2)$ in $t$.  At least one of the first
$(|\mathcal R|-1)(m-2)+1$ positive integers is therefore not a root
over this field.  Use the first such integer.  This is a deterministic
choice, gives integer polynomial coordinates, and has every asserted
nonzero response.  Applying the same construction to $A_*$ gives
$w_0$.

Include $D$, all nonzero response polynomials, and the required row
minors in the finite polynomial list used to choose $M$ and the gap
scales.  There is no dependence on the functions in this list.
Delete tuples with a zero polynomial value or a repeated prime entry,
and tuples for which $\pi^{e_0}\mid D$ for some $\pi\le w$.
The deleted probability is $o(1)$; the first two parts even have
superpolynomially small probability in $V$.
Deletion at this point remains legitimate for weighted expectations:
for any product of distinct row weights its integral over a prime-only
domain $E$ equals $2^t\PP(E)+o(1)$, where $t$ is the number of weights,
by expansion and the restricted-domain assertion of
Proposition~\ref{prop:linear-forms}.  A domain of probability $o(1)$
therefore has $o(1)$ weighted mass.  We henceforth normalize the prime
law on the remaining good tuples, whose probability tends to one.

For integrality, write $D=\epsilon S Q$, where $\epsilon\in\{-1,1\}$,
$S=|D|_{\le w}$, and $Q=|D|_{>w}$.  On the good tuples
$S\mid W^{e_0-1}$.  Since $M/c_k$ is an integer divisible by
$W^{e_0+1}$, the expression in \eqref{eq:correlation-integer-directions}
can be written
\[
 v_{R,k}=
 \epsilon\frac{M/c_k}{S}\,c_{a_*}w_{R,k}\frac{D}{d_R}.
\]
All its factors are integers, and $(M/c_k)/S$ is divisible by $W$.
The assertion for $v_0$ follows directly from its formula.
The degree and coefficient bounds for the polynomial vectors give
$|v_{R,k}|+|v_{0,k}|\le(P_l^++V)^{O_m(1)}$; all the scales $c_k$
are at most $M\le V$.

Direct substitution proves \eqref{eq:correlation-direction-responses}.
More generally,
\[
 \ell_I(v_R)=M(\boldsymbol p)\frac{c_{a_*}}{c_{a(I)}}
                   \frac{A_Iw_R}{d_R},\qquad
 \ell_I(v_0)=M\frac{c_{a_*}}{c_{a(I)}}A_Iw_0.
\]
Outside primes dividing $D(A_Iw_R)$, respectively $A_Iw_0$, the first,
respectively second, expression is a unit at $w<\pi\le V$.
The smooth factors are units there.  This proves both the response
claim and the finite-polynomial nature of all exceptional tests.
\end{proof}

\subsection{Additive elimination and removal of the remaining weights}

\begin{lemma}[Weighted additive elimination]
\label{lem:additive-elimination}
With the rows, good-tuple law, and modulus of
Lemma~\ref{lem:row-directions}, fix $J_0\in\N$ and put
\begin{equation}
 L(\boldsymbol p)=\left\lfloor\frac{R}{J_0M(\boldsymbol p)}\right\rfloor.
 \label{eq:correlation-shift-length}
\end{equation}
Independently conditional on the primes, let every $u_R^0,u_R^1$ be
uniform on $[0,L(\boldsymbol p))\cap\Z$, for $R\ne *$.
Then
\begin{equation}
 \left|\E_{\boldsymbol p,z}\prod_{I\in\mathcal R}f_I(\ell_I(z))
 \right|^{2^d}
 \le C_m\left|\E_{\boldsymbol p,z,u}
       \prod_{\omega\in\{0,1\}^d}
       g_*\left(\ell_*(z)+M(\boldsymbol p)
                    \sum_{R\ne *}u_R^{\omega_R}\right)\right|+o(1).
 \label{eq:additive-elimination-output}
\end{equation}
The constants are independent of $J_0$; the error is for fixed $J_0$.
\end{lemma}

\begin{proof}
Weighted Cauchy--Schwarz will leave a cube of $g_*$ multiplied by
retained row weights.  We first form that cube, then use the direction
$v_0$ to remove the retained weights in a weighted second moment.

Every interval in \eqref{eq:correlation-shift-length} has length
dominating all fixed powers of $V$, uniformly over the prime tuples.
Indeed $M(\boldsymbol p)$ is bounded by a fixed power of $P_l^++V$,
whereas $R$ dominates all such powers.  Insert independent uniform
$u_R$ on these intervals and translate
\[
 z\longmapsto z+\sum_{R\ne *}v_Ru_R.
\]
The translations are in $W\Z^m$.  Their sizes are at most
$R(P_l^++V)^{O_m(1)}$, and hence are negligible at every pivot cutoff,
even after multiplying the total-variation errors by $V^{O_m(1)}$.
Lemma~\ref{lem:sampling} therefore changes the average by $o(1)$.

Eliminate the nondistinguished rows in a fixed order.  At the step
assigned to row $R$, all its current copies are independent of $u_R$,
because $\ell_R(v_R)=0$.  Let $H_{\mathrm{out}}$ be their product,
and let $\Omega_R$ be the product of their bounds $W_R$.
All other factors, including weights retained at earlier steps, belong
to $H_{\mathrm{in}}$.  With all variables other than $u_R$ outside,
the exact inequality is
\begin{align}
 &\left|\E_{\mathrm{out}}H_{\mathrm{out}}
                   \E_{u_R}H_{\mathrm{in}}\right|^2\notag\\
 &\hspace{6mm}\le
   (\E_{\mathrm{out}}\Omega_R)
   \E_{\mathrm{out},u_R^0,u_R^1}
       \Omega_R H_{\mathrm{in}}(u_R^0)H_{\mathrm{in}}(u_R^1).
 \label{eq:additive-weighted-cs}
\end{align}
It follows by applying Cauchy--Schwarz with density $\Omega_R$ and
using $|H_{\mathrm{out}}|\le\Omega_R$.  In particular, the assigned
functions disappear and their weights occur once, while every factor
inside the average is duplicated.  All prime parameters remain
outside throughout this procedure.

We verify every use of Proposition~\ref{prop:linear-forms} in
\eqref{eq:additive-weighted-cs}.  A current copy of a base row $I$ has
the form
\[
 \ell_I(z)+\sum_Q\ell_I(v_Q)u_Q^{\eta_Q},
\]
where $\eta_Q$ is one of the already duplicated choices or the single
current variable.  Distinct base rows are distinguished by their
$z$-coefficient vectors.  For two copies of the same row, some
duplicated coordinate $Q\ne I$ has different choices, and its response
$\ell_I(v_Q)$ is nonzero.  Their difference therefore has a nonzero
coefficient on an independent variable $u_Q^0$ or $u_Q^1$.
Taking a minor with an anchor $z$-column proves pairwise independence
at every relevant prime outside the response-polynomial tests of
Lemma~\ref{lem:row-directions}.  A row's anchor coefficient is still
a monomial of pool primes, so each row is primitive at all
$w<\pi\le V$.  Finally, conditional on the primes the base variables
are independent harmonic variables or independent interval variables
whose lengths dominate powers of $V$.  Their joint residues modulo
the product of any fixed number of sampled divisors are uniform with
error smaller than every fixed negative power of $V$.
These facts check primitivity, pair independence with polynomial
exceptions, and the sampling hypotheses.  Expanding the weights now
shows that every prefactor in \eqref{eq:additive-weighted-cs} is bounded
by a constant depending only on the number of row copies.

After all $d$ steps, the nondistinguished functions have disappeared.
The remaining target factor is
\[
 G(z,u)=\prod_{\omega\in\{0,1\}^d}
 g_*\left(\ell_*(z)+M(\boldsymbol p)
                   \sum_Ru_R^{\omega_R}\right).
\]
For each eliminated row $I$, there is one retained weight for each
choice of the duplicated variables in the directions $R\ne I$.
Thus the retained product is
\begin{equation}
 \Psi(z,u)=
 \prod_{I\ne *}\ \prod_{\eta\in\{0,1\}^{[d]\setminus\{I\}}}
 W_I\left(\ell_I(z)+\sum_{R\ne I}\ell_I(v_R)u_R^{\eta_R}\right),
 \label{eq:correlation-retained-weights}
\end{equation}
after identifying the nondistinguished rows with $[d]$.
There are $t=d2^{d-1}$ weight factors.  Iteration has proved the left
side of \eqref{eq:additive-elimination-output} is bounded by a constant
times $|\E G\Psi|+o(1)$.

It remains to justify replacing $\Psi$ by $2^t$ in this correlation;
its mean alone would not justify that replacement.  Translate $z$
once more by $v_0u_0$, with $u_0$ uniform on $[0,R)\cap\Z$.
The sampling error is again $o(1)$, and the target $G$ is unchanged
because $\ell_*(v_0)=0$.  Put
\[
 B(z,u)=\prod_{\omega\in\{0,1\}^d}
  \left(1+\nu_{a_*}\left(\ell_*(z)+M(\boldsymbol p)
                 \sum_Ru_R^{\omega_R}\right)\right),
 \qquad H(z,u)=\E_{u_0}\Psi(z+v_0u_0,u).
\]
Then $|G|\le B$.  The same row check gives
\begin{equation}
 \E B=2^{2^d}+o(1),\qquad
 \E BH=2^{2^d+t}+o(1),\qquad
 \E BH^2=2^{2^d+2t}+o(1).
 \label{eq:auxiliary-weight-moments}
\end{equation}
For completeness, in the last expression duplicate $u_0$ independently
as $u_0^0,u_0^1$.  Copies from distinct base rows are distinguished on
$z$; copies from one base row and different old shift choices are
distinguished on an old $u_R^\eta$; copies with identical old choices
but different $u_0$ branches are distinguished by the nonzero response
$\ell_I(v_0)$.  Target rows have distinct cube shift choices and no
$u_0$ response.  All forms in every expanded moment are thus distinct
and satisfy exactly the primitivity and polynomial-minor conditions
checked above.  Expanding each $1+\nu$ and replacing every product of
$\nu$'s by its main term $1$ gives all three powers of $2$ in
\eqref{eq:auxiliary-weight-moments}.

Consequently
\[
 \E B(H-2^t)^2=o(1).
\]
Cauchy--Schwarz with density $B$ now yields
\[
 |\E G(H-2^t)|
 \le (\E B)^{1/2}\bigl(\E B(H-2^t)^2\bigr)^{1/2}=o(1).
\]
This is the required weighted replacement.  Combining it with the
$d$ inequalities \eqref{eq:additive-weighted-cs} proves
\eqref{eq:additive-elimination-output}.
\end{proof}

\subsection{A one-variable cube estimate}

\begin{proposition}[Uniform correlation test]\label{prop:correlation-test}
For every $m\ge2$ and every nonsingleton target support $J_*$, the
construction above determines a finite family of integer polynomial
templates $D$, with good prime-tuple laws, such that the following holds
for every fixed master chain and valid gap $l$.
There are an integer $1\le d\le K_m-1$, a constant $C_m$, and
$\theta>0$, all bounded in terms of $m$ alone, for which
\begin{equation}
 \boxed{\quad
 |\mathcal C|
 \le o(1)+C_m\left|
  \E_{\boldsymbol p,y,u}
  \prod_{\omega\subseteq[d]}
  g_*\left(y+M(\boldsymbol p)
                  \sum_{R\in\omega}(u_R^1-u_R^0)\right)
                       \right|^{\theta}.
 \quad}
 \label{eq:correlation-test}
\end{equation}
Here $y$ has law $\mu_{i_{a_*}}$, independently of the primes and
shifts; the primes have the normalized good-tuple law of
Lemma~\ref{lem:row-directions}; $M(\boldsymbol p)$ is given by
\eqref{eq:correlation-modulus}; and the shifts have the conditional laws
in \eqref{eq:correlation-shift-length}.  One may take
$\theta=2^{-(q_{\mathrm{mask}}+d)}$.  The gap length $R_l$ is divisible by every
$M(\boldsymbol p)$ that occurs.

The estimate is uniform over all the masks and functions in
\eqref{eq:correlation-initial}, including $g_*$, and over each fixed
finite set of positive integers $J_0$.  Neither $d$, $C_m$, nor
$\theta^{-1}$ depends on $N$, on the scales, or on the functions.
The polynomial templates and direction choices are fixed before $M$
and the prime pools are chosen.  In particular, the same tests can be
used both for a bounded error and for an error bounded by
$1+\nu_{a_*}$.
\end{proposition}

\begin{proof}
Only the pushforward of the cube root remains after
Lemmas~\ref{lem:mask-removal} and \ref{lem:additive-elimination}.
For fixed primes and shifts, write
\begin{equation}
 Y_* = \ell_*(z)+M(\boldsymbol p)\sum_Ru_R^0
      = kz_{a_*}+h,
 \qquad k=A_{*,a_*}.
 \label{eq:correlation-root}
\end{equation}
The target support remains $J_*$, so it contains some $j<a_*$.
The coefficient
$b=(c_j/c_{a_*})A_{*,j}$ is coprime to $k$:
the prime slots in different columns are disjoint, their numerical
entries are distinct on the good tuples, and all scale ratios are
$w$-smooth whereas all pool primes exceed $V\ge w$.
Also $k$ is coprime to $W$ and $W\mid h$.

Condition first on every $z$-variable except $z_{a_*}$.  Put
$X=X_{i_{a_*}}$ and $\delta_W=\phi(W)/W$.  Uniformly in the variables
being conditioned on, one may take
\[
 0\le h\le H:=M(P_l^+)^{B_m}\sum_{j<a_*}X_{i_j}^2+dR/J_0
\]
for a fixed $B_m$.  The logarithm of $X$ dominates every fixed power
of $H+k+W$, and in particular $H<X/2$ eventually.
The dilation and translation calculation of Lemma~\ref{lem:sampling}
gives
\begin{equation}
 \mathcal L(kz_{a_*}+h)
   = k\1_{y\equiv h\pmod{k}}\mu_{i_{a_*}}(dy)+\mathcal E_h,
 \qquad \|\mathcal E_h\|_{\mathrm{TV}}=o(V^{-C})
 \label{eq:correlation-root-progression}
\end{equation}
for every fixed $C$.  More explicitly, the total-mass error is at most
\begin{equation}
 O\left(\frac{\log(2k)}{\log X}+\frac H X+
                  \frac{Wk^2}{\delta_W X\log X}\right).
 \label{eq:correlation-root-tv-bound}
\end{equation}
To see this directly, let $Z_i$ be the normalizing constant of $\mu_i$,
so $Z_i\asymp\delta_W\log X$.  The exact image density is
$k/(Z_i(y-h))$ on $[kX+h,kX^2+h)$ in the class $h\pmod k$,
with $(y,W)=1$.  The comparison density is $k/(Z_i y)$ on
$[X,X^2)$ in the same class.  The $W$-unit restriction agrees because
$W\mid h$ and $(k,W)=1$.  On the intersection of the intervals the
relative discrepancy is at most $H/X$.  Harmonic progression counting
on the two remaining boundary intervals gives a logarithmic main
term $O(\delta_W\log(2k))$ and error $O(Wk^2/X)$ before normalization.
This proves \eqref{eq:correlation-root-tv-bound}.  All three terms
are smaller than every fixed negative power of $V$ by scale separation;
the crude bound $\delta_W^{-1}\le W$ already suffices here.

Now condition on all remaining variables except $z_j$.
Since $i_j>l$, Lemma~\ref{lem:sampling} makes $z_j$ uniform modulo $k$;
its relative error is at most
\[
 O\left(\frac{Wk^2}{\delta_W X_{i_j}\log X_{i_j}}\right)=o(V^{-C})
\]
for every fixed $C$.
The same holds for $h=bz_j+h_0$ modulo $k$, because $(b,k)=1$.
Averaging the density in \eqref{eq:correlation-root-progression}
therefore gives
\[
 \E_{z_j} k\1_{y\equiv h\pmod{k}}=1+o(V^{-C})
\]
uniformly in $y$ and in the conditioned primes and shifts.
Thus the conditional law of $Y_*$, after integrating all $z$'s, differs
from $\mu_{i_{a_*}}$ by $o(V^{-C})$ in total mass, uniformly in the
primes and shifts.  This argument also covers $k=1$, when there is no
residue condition to average.

Every remaining cube argument can be expressed using this root as
\[
 Y_*+M(\boldsymbol p)
                 \sum_{R\in\omega}(u_R^1-u_R^0).
\]
There is no other $z$-dependent factor left.  The product of target
bounds is at most $(1+V)^{2^d}$, so the total-mass error just proved is
still $o(1)$ in the cube expectation.  Substitution in
\eqref{eq:additive-elimination-output} and then
\eqref{eq:mask-removal-output} gives
\eqref{eq:correlation-test} with the stated exponent.  All constants
arise from fixed counts of row weights and squares, hence depend only
on $m$.  At least one other row exists, so $d\ge1$.

Finally, $M(\boldsymbol p)$ divides $M|D(\boldsymbol p)|$, whose
divisibility into $R_l$ was arranged in
Lemma~\ref{lem:master-scales}.  The finite template assertion follows
from the deterministic mask and kernel-vector choices.  Every
sampling estimate and every restricted linear-forms estimate used
above was uniform over the functions; taking the maximum of the errors
over finitely many $J_0$'s proves the last uniformity assertion.
\end{proof}

\begin{remark}[Tests used in the next section]
\label{rem:correlation-dual-interface}
The base vertex in \eqref{eq:correlation-test} is $g_*(y)$.
All its other vertices may therefore be regarded as inputs in a dual
test
\[
 \mathcal D(y)=\E_{\boldsymbol p,u}e(\boldsymbol p)
  \prod_{\emptyset\ne\omega\subseteq[d]}
  h_{\omega,\boldsymbol p}\left(y+M(\boldsymbol p)
              \sum_{R\in\omega}(u_R^1-u_R^0)\right),
 \quad |e|\le1,\quad |h_{\omega,\boldsymbol p}|\le1+\nu_{a_*}.
\]
In particular, choosing every $h_{\omega,\boldsymbol p}=g_*$ and
$e=1$ recovers its cube expectation as $\E_{\mu_{i_{a_*}}}g_*\mathcal D$.
The enlarged test class, including independently chosen inputs and
bounded prime-dependent signs, is fixed without reference to a dense
model.  Proposition~\ref{prop:correlation-test} itself asserts the
inequality for the repeated target function; the pseudorandomness of
the enlarged test class is proved in Section~\ref{sec:prediction}.
\end{remark}
